\documentclass[11pt,a4paper]{article}
% ============================================================
% Shared preamble for Claude Experiments reports (pdflatex)
% Include this AFTER \documentclass in each report:
%   \documentclass[11pt,a4paper]{article}
%   % ============================================================
% Shared preamble for Claude Experiments reports (pdflatex)
% Include this AFTER \documentclass in each report:
%   \documentclass[11pt,a4paper]{article}
%   % ============================================================
% Shared preamble for Claude Experiments reports (pdflatex)
% Include this AFTER \documentclass in each report:
%   \documentclass[11pt,a4paper]{article}
%   \input{../preamble}
% ============================================================
\usepackage[T1]{fontenc}
\usepackage[utf8]{inputenc}
\usepackage[english]{babel}
\usepackage[a4paper,top=2.5cm,bottom=2.5cm,left=2.5cm,right=2.5cm]{geometry}
\usepackage{amsmath,amssymb}
\usepackage{graphicx}
\usepackage{subcaption}
\usepackage{booktabs}
\usepackage{xcolor}
\usepackage[colorlinks=true,allcolors=blue]{hyperref}
\usepackage{microtype}
\usepackage{float}
\usepackage{array}
\microtypesetup{expansion=false}

% Color scheme
\definecolor{colMicro}{HTML}{222222}
\definecolor{colExact}{HTML}{2878b5}
\definecolor{colPlacement}{HTML}{e07a1f}
\definecolor{colBoundary}{HTML}{429c74}
\definecolor{colCurrent}{HTML}{c23b55}
\definecolor{colMM}{HTML}{c23b55}

\providecommand{\SIG}{\textcolor{green!50!black}{\textbf{*}}}
\providecommand{\NS}{\textcolor{gray}{n.s.}}
\providecommand{\DRMS}{\Delta\mathrm{RMST}}
\providecommand{\ci}[2]{$[{#1},\;{#2}]$}

% ============================================================
\usepackage[T1]{fontenc}
\usepackage[utf8]{inputenc}
\usepackage[english]{babel}
\usepackage[a4paper,top=2.5cm,bottom=2.5cm,left=2.5cm,right=2.5cm]{geometry}
\usepackage{amsmath,amssymb}
\usepackage{graphicx}
\usepackage{subcaption}
\usepackage{booktabs}
\usepackage{xcolor}
\usepackage[colorlinks=true,allcolors=blue]{hyperref}
\usepackage{microtype}
\usepackage{float}
\usepackage{array}
\microtypesetup{expansion=false}

% Color scheme
\definecolor{colMicro}{HTML}{222222}
\definecolor{colExact}{HTML}{2878b5}
\definecolor{colPlacement}{HTML}{e07a1f}
\definecolor{colBoundary}{HTML}{429c74}
\definecolor{colCurrent}{HTML}{c23b55}
\definecolor{colMM}{HTML}{c23b55}

\providecommand{\SIG}{\textcolor{green!50!black}{\textbf{*}}}
\providecommand{\NS}{\textcolor{gray}{n.s.}}
\providecommand{\DRMS}{\Delta\mathrm{RMST}}
\providecommand{\ci}[2]{$[{#1},\;{#2}]$}

% ============================================================
\usepackage[T1]{fontenc}
\usepackage[utf8]{inputenc}
\usepackage[english]{babel}
\usepackage[a4paper,top=2.5cm,bottom=2.5cm,left=2.5cm,right=2.5cm]{geometry}
\usepackage{amsmath,amssymb}
\usepackage{graphicx}
\usepackage{subcaption}
\usepackage{booktabs}
\usepackage{xcolor}
\usepackage[colorlinks=true,allcolors=blue]{hyperref}
\usepackage{microtype}
\usepackage{float}
\usepackage{array}
\microtypesetup{expansion=false}

% Color scheme
\definecolor{colMicro}{HTML}{222222}
\definecolor{colExact}{HTML}{2878b5}
\definecolor{colPlacement}{HTML}{e07a1f}
\definecolor{colBoundary}{HTML}{429c74}
\definecolor{colCurrent}{HTML}{c23b55}
\definecolor{colMM}{HTML}{c23b55}

\providecommand{\SIG}{\textcolor{green!50!black}{\textbf{*}}}
\providecommand{\NS}{\textcolor{gray}{n.s.}}
\providecommand{\DRMS}{\Delta\mathrm{RMST}}
\providecommand{\ci}[2]{$[{#1},\;{#2}]$}


\title{CE-03: MicroMacro Error vs.\ Beta Ratio\\
       \large Experiment Report --- Two-Scale CTMC Project}
\author{Claude Experiments Pipeline}
\date{2026-07-31}

\begin{document}
\maketitle

\begin{abstract}
This report quantifies how the ratio $r = \beta_{\text{out}}/\beta_{\text{in}}$
between inter- and intra-community infection rates affects the MicroMacro
approximation error.  Five scenarios sweep $r \in \{0.1, 0.3, 1.0, 2.0, 5.0\}$
on a fixed chain-of-cliques topology ($L=10$, $n=100$, $d=6$, $k=5$)
with 100 paired replications each.
The primary finding is a \textbf{sign reversal} in $\DRMS$:
MicroMacro arrives \emph{later} than full Gillespie at low $r$
($\DRMS = {+}4796$ at $r=0.1$), and \emph{earlier} at high $r$
($\DRMS = -1355$ at $r=1.0$, $-1638$ at $r=2.0$).
The crossover occurs between $r=0.3$ and $r=1.0$.
At $r \ge 2$ the full-Micro epidemic rarely reaches the final community
within the simulation horizon (completion 54\% and 1\% respectively),
while MicroMacro completes in all runs---indicating a qualitative
false-positive bias under strong inter-community coupling.
\textit{Ukrainian summary}: see \texttt{docs/ce03\_beta\_ratio/report.md}.
\end{abstract}

\section{Experimental Design}

\subsection{Objective}
Determine how the ratio of inter-community to intra-community infection rate
modulates MicroMacro approximation error, and characterise the direction
(premature vs.\ delayed arrival) across the parameter range.

\subsection{Network and Parameters}
\begin{itemize}
  \item Topology: chain-of-cliques, $L=10$ communities, $n=100$ nodes each,
        degree $d=6$, $k=5$ inter-community edges per boundary.
  \item Intra-community infection rate: $\beta_{\text{in}} = 0.003$.
  \item Inter-community infection rate: $\beta_{\text{out}} = r \cdot \beta_{\text{in}}$.
  \item Rate ratios: $r \in \{0.1,\; 0.3,\; 1.0,\; 2.0,\; 5.0\}$.
  \item Recovery: $\gamma = 0$ (pure SI model).
  \item Replications: 100 paired runs per scenario (same network and seed).
  \item Horizon $T_{\text{end}}$: scenario-dependent (20{,}000 to 2{,}000;
        shorter for larger $r$ to keep runtimes feasible).
\end{itemize}

\subsection{Primary Metric}
\[
  \DRMS = \overline{T_{\text{MM}}} - \overline{T_{\text{Micro}}}
\]
where arrival times use restricted-mean estimates with censoring at
$T_{\text{end}}$ (censored runs contribute $T_{\text{end}}$ to the mean).
A positive $\DRMS$ indicates MicroMacro predicts \emph{delayed} arrival;
negative indicates \emph{premature} arrival.
Bootstrap 95\% confidence intervals use 10{,}000 resamples.
Relative Wasserstein distance $\text{rel-}W_1$ measures distributional
discrepancy normalised by the Micro mean arrival time.

\section{Results}

\begin{table}[H]
\centering
\caption{CE-03 results: MicroMacro error by beta ratio.
  $\DRMS$ = censored-mean difference (MicroMacro $-$ Micro).
  $r = \beta_{\text{out}}/\beta_{\text{in}}$.
  ``Comp.'' = fraction of Micro runs reaching the final community
  by $T_{\text{end}}$ (MicroMacro: 1.00 for all scenarios).
  All CIs exclude zero.}
\label{tab:ce03-main}
\begin{tabular}{rrrrrrc}
\toprule
$r$ & $\beta_{\text{out}}$ & $\DRMS$ & \multicolumn{2}{c}{95\% CI} & rel-$W_1$ & Micro comp. \\
\cmidrule(lr){4-5}
 & & & lower & upper & & \\
\midrule
0.1 & 0.0003 & $+4796$ & $+4384$ & $+5213$ & 1.624 & 1.00 \\
0.3 & 0.0009 & $+562$  & $+300$  & $+812$  & 0.297 & 1.00 \\
1.0 & 0.0030 & $-1355$ & $-1494$ & $-1220$ & 0.456 & 1.00 \\
2.0 & 0.0060 & $-1638$ & $-1746$ & $-1527$ & 0.558 & 0.54 \\
5.0 & 0.0150 & $-1213$ & $-1283$ & $-1141$ & 0.566 & 0.01 \\
\bottomrule
\end{tabular}
\end{table}

\subsection{Key Observations}

\begin{enumerate}
  \item \textbf{Sign reversal.}  $\DRMS$ is positive at $r \le 0.3$ and
        negative at $r \ge 1.0$.  The crossover $r^*$ lies between 0.3 and 1.0.
  \item \textbf{Large magnitude at extremes.}  $|\DRMS| \approx 4800$ at $r=0.1$
        and $\approx 1600$ at $r=2.0$---both large relative to the horizon.
  \item \textbf{Distributional divergence.}  $\text{rel-}W_1 = 1.624$ at $r=0.1$
        (distributions differ by more than the Micro mean arrival time).
  \item \textbf{Completion collapse at high $r$.}  At $r=2.0$ only 54\% of
        full-Micro runs reach the final community within $T_{\text{end}}=3000$;
        at $r=5.0$ only 1\%.  MicroMacro always reports full traversal,
        yielding a 99 percentage-point false-positive epidemic rate at $r=5.0$.
\end{enumerate}

\section{Interpretation}

\subsection{Mechanism at Low $r$ (delayed arrival, $\DRMS > 0$)}

When $\beta_{\text{out}} \ll \beta_{\text{in}}$, intra-community spread is fast
and inter-community spread is rare.  In the full Gillespie (Micro) model,
the $k=5$ gateway nodes on the boundary of the source community can fire their
cross-community edges as soon as \emph{they themselves} become infected,
regardless of the community-wide infected fraction $\bar{i}(t)$.

MicroMacro uses the aggregate import hazard
$\lambda_{\text{macro}}(t) = \beta_{\text{out}} \cdot \bar{i}(t)$,
proportional to the community-wide mean.  When $\beta_{\text{out}}$ is
tiny, this hazard remains negligible until $\bar{i}(t)$ grows substantially,
so the macro import event fires \emph{late}.  The direct gateway edges in Micro
can fire earlier, giving $\DRMS > 0$.

\subsection{Mechanism at High $r$ (premature arrival, $\DRMS < 0$)}

When $\beta_{\text{out}} \ge \beta_{\text{in}}$, the aggregate hazard
$\lambda_{\text{macro}}(t) = \beta_{\text{out}} \cdot \bar{i}(t)$ becomes
large as soon as any infecteds appear in the source community.  MicroMacro
fires imports aggressively and places each at a uniformly random node in the
target community, artificially accelerating chain traversal.

In full Micro, the epidemic must propagate through only $k=5$ physical edges
per boundary.  Even with large $\beta_{\text{out}}$, traversing nine
consecutive boundaries in a chain is constrained by within-community buildup
at each step.  Micro rarely completes the chain within a short $T_{\text{end}}$
(1\% at $r=5.0$), while MicroMacro always does.

\subsection{Comparison with CE-01 Baseline}

CE-01 used homogeneous rates ($r=1$) and found $\DRMS = -209$ with
rate abstraction as the dominant error source.  CE-03 at $r=1.0$ gives
$\DRMS = -1355$; the larger magnitude is consistent with a shorter
$T_{\text{end}} = 4000$ (vs.\ 5000 in CE-01).  Both confirm premature
arrival at $r=1$.

\section{Conclusion}

CE-03 reveals a qualitative change in MicroMacro error direction as a function
of $r$:

\begin{itemize}
  \item \textbf{Low $r$ ($r < r^* \approx 0.5$--$0.8$):}
        MicroMacro predicts \emph{delayed} arrival.  The aggregate hazard
        underweights early gateway activity, causing the macro import to fire too late.
  \item \textbf{High $r$ ($r > r^*$):}
        MicroMacro predicts \emph{premature} arrival.  The aggregate hazard
        overestimates the effective multi-boundary crossing rate.
        At $r=5.0$ the qualitative prediction is wrong: Micro rarely spreads
        across the full chain, while MicroMacro always does (99\% false-positive rate).
  \item \textbf{Crossover near $r^* \approx 0.5$--$0.8$:}
        Both error directions cancel.  This cancellation is topology-dependent,
        not a regime of correct approximation.
\end{itemize}

A single scalar correction to $\beta_{\text{out}}$ cannot remove the bias
across the full parameter range.  Structural corrections---such as using the
actual gateway-node infected count rather than the community mean---are needed.

\section*{Data Location}
\texttt{output/claude\_experiments/ce03\_beta\_ratio/summary.csv}

\section*{Related Experiments}
CE-01 (ablation: rate abstraction vs.\ placement), CE-04 (SIR: recovery effect).
\end{document}
